% Public manuscript source.
\documentclass[11pt]{article}

\usepackage[T1]{fontenc}
\usepackage{lmodern}
\usepackage[margin=1.05in]{geometry}
\usepackage{amsmath,amsthm}
\usepackage{enumitem}
\usepackage{microtype}
\usepackage{xcolor}
\usepackage[colorlinks=true,linkcolor=blue!55!black,citecolor=blue!55!black,urlcolor=blue!55!black]{hyperref}

\newtheorem{theorem}{Theorem}[section]
\newtheorem{lemma}[theorem]{Lemma}
\newtheorem{proposition}[theorem]{Proposition}
\newtheorem{corollary}[theorem]{Corollary}
\theoremstyle{definition}
\newtheorem{definition}[theorem]{Definition}

\newcommand{\dd}{\delta}

\title{An 18-Vertex Counterexample to the
Oriented Berge--Fulkerson Conjecture}
\author{\textbf{Xiaobo Wu}\\
  \textit{Faculty of Science, National University of Singapore}\\
  \texttt{xiaobo.wu@u.nus.edu}}
\date{\small\href{https://doi.org/10.5281/zenodo.22282127}{DOI: 10.5281/zenodo.22282127}}

\begin{document}
\maketitle

\begin{abstract}
The oriented Berge--Fulkerson conjecture (o6c4c) asserts that every finite
loopless bridgeless graph admits six oriented cycles such that every edge
occurs in exactly four indexed members, twice in each direction.  We construct
a simple cubic $3$-edge-connected counterexample on $18$ vertices, obtained
from the Petersen graph by replacing the four vertices of a closed
neighbourhood by triangles.  Any putative cover forces the six complementary
perfect matchings to be the unique lifts of the six perfect matchings of the
Petersen graph.  A fixed integer linear combination of five edge-balance
equations then yields $4z_0=0$ for a sign $z_0\in\{-1,1\}$, a contradiction.
The proof is self-contained; exhaustive computation is used only as an
independent verification.
\end{abstract}

\section{Introduction}

Ulyanov~\cite[Conjecture~1.1]{Ulyanov} proposed the oriented
Berge--Fulkerson conjecture, abbreviated o6c4c, as an oriented strengthening
of the $6$-cycle $4$-cover reformulation of the Berge--Fulkerson
conjecture~\cite{Fulkerson}.  Following Ulyanov's terminology and the standard
circuit-cover convention~\cite{Zhang}, a cycle is an even subgraph and need
not be connected.

For the cubic graph considered below, every cycle is a disjoint union of
circuits.  We orient such a cycle componentwise, assigning each circuit one
of its two cyclic orientations.  Equivalently, the resulting directed
subgraph has equal indegree and outdegree at every vertex.  Thus no convention
for orienting an even subgraph at a vertex of degree greater than two enters
the proof.

\begin{theorem}\label{thm:main}
There exists a finite simple cubic $3$-edge-connected graph $H$ on $18$
vertices that admits no oriented $6$-cycle $4$-cover.
\end{theorem}

This does not conflict with the computational verification for snarks on at
most $30$ vertices reported in~\cite[Section~2.3]{Ulyanov}: the graph $H$
constructed here has girth $3$ and cyclic edge-connectivity $3$, and hence is
outside the class considered there.

The graph is a four-triangle expansion of the Petersen graph.  The argument
has two rigid stages.  First, local degree counting turns every putative cover
into a perfect-matching double cover, and the matching structure of the
Petersen graph makes that cover unique up to permuting its six indexed
members.  Second, direction balance on five selected edges gives an
inconsistent integer system.

\section{Definitions and the signed formulation}

All graphs are finite and loopless.  The counterexample is simple, so no
convention concerning parallel edges affects the result.

\begin{definition}\label{def:o6c4c}
A \emph{cycle} in a graph $G$ is an edge set $F\subseteq E(G)$ such that
$d_F(v)$ is even for every $v\in V(G)$; it need not be connected.  An
orientation of $F$ is \emph{Eulerian} if
\[
d_F^+(v)=d_F^-(v)\qquad(v\in V(G)).
\]

An \emph{oriented $6$-cycle $4$-cover}, or \emph{o6c4c}, is an indexed
six-tuple $(F_0,\ldots,F_5)$ of oriented cycles such that every edge occurs in
exactly four members and, among those four occurrences, is directed twice in
each direction.  Repetitions among the underlying edge sets are permitted.
\end{definition}

When $G$ is cubic, every cycle has degree $0$ or $2$ at each vertex.  Hence
its nontrivial connected components are circuits, and an Eulerian orientation
is precisely a choice of one of the two cyclic orientations on each circuit
component.

Fix a reference direction for every edge of $G$.  An oriented cycle $F_i$ is
encoded by a function
\[
s_i\colon E(G)\longrightarrow\{-1,0,1\},
\]
where $s_i(e)=0$ if $e\notin F_i$, $s_i(e)=1$ if the orientation of $e$ in
$F_i$ agrees with the reference direction, and $s_i(e)=-1$ otherwise.
Let $\varepsilon(v,e)=1$ if the reference direction of $e$ points away from
$v$, let $\varepsilon(v,e)=-1$ if it points towards $v$, and let
$\varepsilon(v,e)=0$ if $e$ is not incident with $v$.

\begin{lemma}[Signed formulation]\label{lem:signed}
A graph $G$ has an oriented $6$-cycle $4$-cover if and only if there are
functions
\[
s_0,\ldots,s_5\colon E(G)\longrightarrow\{-1,0,1\}
\]
satisfying
\begin{align}
\sum_{e\in E(G)}\varepsilon(v,e)s_i(e)&=0
    &&(v\in V(G),\ 0\leq i\leq5), \label{eq:flow}\\
\sum_{i=0}^{5}|s_i(e)|&=4
    &&(e\in E(G)), \label{eq:support}\\
\sum_{i=0}^{5}s_i(e)&=0
    &&(e\in E(G)). \label{eq:balance}
\end{align}
\end{lemma}

\begin{proof}
Equation~\eqref{eq:flow} is precisely vertex balance in each member, while
\eqref{eq:support} says that four indexed members contain the edge.  Four
nonzero signs in $\{-1,1\}$ have sum zero exactly when two are positive and two
are negative, giving~\eqref{eq:balance}.  The converse reconstruction is
immediate.

Reversing the reference direction of an edge negates both
$\varepsilon(\,\cdot\,,e)$ and all six values $s_i(e)$.  Thus
\eqref{eq:flow} is unchanged, while \eqref{eq:support} and
\eqref{eq:balance} remain equivalent to their original forms.
\end{proof}

\section{The 18-vertex graph}

Let $P$ be the Petersen graph on vertices $0,\ldots,9$ with edge set
\begin{equation}\label{eq:petersen-edges}
E(P)=\{01,04,05,12,16,23,27,34,38,49,57,58,68,69,79\}.
\end{equation}
Here and below, $vw$ denotes the edge $\{v,w\}$.  Set
$X=N_P[0]=\{0,1,4,5\}$.  For each $v\in X$, replace $v$ by a triangle whose
three vertices are ports $p_{vw}$ indexed by the neighbours $w$ of $v$.
Join the three ports belonging to the same expanded vertex in a triangle, and
denote an unexpanded vertex $v$ by $u_v$.  For an old edge $vw\in E(P)$, its
image in $H$ is
\[
\begin{cases}
p_{vw}p_{wv}, & v,w\in X,\\
p_{vw}u_w,    & v\in X,\ w\notin X,\\
u_vp_{wv},    & v\notin X,\ w\in X,\\
u_vu_w,       & v,w\notin X.
\end{cases}
\]
Denote the resulting graph by $H$.

\begin{lemma}\label{lem:expansion}
Replacing a vertex of a cubic $3$-edge-connected graph by a triangle, with its
three old incident edges attached to distinct triangle vertices, preserves
$3$-edge-connectivity.
\end{lemma}

\begin{proof}
Suppose the expanded graph has an edge cut of size at most two.  If the cut
does not split the new triangle, contracting that triangle gives a cut of the
same size in the original graph.  If the cut splits the triangle, two internal
triangle edges already cross it.  No other edge can cross.  Each external
neighbour is therefore on the same side as its port.  Contract the triangle
and place the old vertex on the side of the lone port.  The two external edges
belonging to the other ports form a cut of size two in the original graph, a
contradiction.
\end{proof}

\begin{proposition}\label{prop:eligible}
The graph $H$ is simple, cubic, $3$-edge-connected, and has $18$ vertices and
$27$ edges.  It has girth $3$ and cyclic edge-connectivity $3$.
\end{proposition}

\begin{proof}
The construction is simple and cubic.  Replacing four vertices adds eight
vertices and twelve edges, so $|V(H)|=10+8=18$ and $|E(H)|=15+12=27$.

For completeness, $P$ is $3$-edge-connected.  It is cubic and triangle-free.
For a nontrivial $S\subset V(P)$, pass to the smaller shore so that
$1\leq |S|\leq5$.  Since
\[
|\dd(S)|=3|S|-2|E(P[S])|,
\]
the triangle-free bounds
\[
|E(P[S])|\leq 0,1,2,4,6
\quad\text{for}\quad |S|=1,2,3,4,5
\]
give $|\dd(S)|\geq3$ in every case.  Apply Lemma~\ref{lem:expansion} four
times.

Each replacement triangle shows that the girth is $3$.  Its three external
edges form a cyclic edge cut: one shore contains the triangle, and the other
shore contains a cycle.  Thus the cyclic edge-connectivity is at most $3$,
and $3$-edge-connectivity gives the reverse inequality.
\end{proof}

An exact numbered edge list for $H$ appears in Appendix~\ref{app:data}.

\section{Rigidity of an unoriented cover}

Suppose that $F_0,\ldots,F_5$ are cycles of $H$ that cover every edge exactly
four times; no orientations are assumed in this section.  Set
\[
N_i=E(H)\setminus F_i
\qquad(0\leq i\leq5).
\]

\begin{lemma}[Support saturation]\label{lem:saturation}
Every $F_i$ is a spanning $2$-factor.  Consequently every $N_i$ is a perfect
matching, and $(N_0,\ldots,N_5)$ covers every edge exactly twice.
\end{lemma}

\begin{proof}
Since $F_i$ is a cycle, $d_{F_i}(v)$ is even for every vertex $v$.  Since
$H$ is cubic,
\[
d_{F_i}(v)\in\{0,2\}.
\]
On the other hand, each of the three edges incident with $v$ belongs to
exactly four of the six supports.  Therefore
\[
\sum_{i=0}^{5}d_{F_i}(v)
=
\sum_{e\ni v}\sum_{i=0}^{5}\mathbf{1}_{e\in F_i}
=
3\cdot4
=
12.
\]
The six summands are each at most $2$, so all of them are equal to $2$.
Thus every $F_i$ is a spanning $2$-factor.

Consequently, $N_i=E(H)\setminus F_i$ meets every vertex exactly once and is
therefore a perfect matching.  Since each edge lies in four of the supports,
it lies in exactly two of the complements.
\end{proof}

The Petersen graph has exactly the following six perfect matchings:
\begin{align}
Q_0&=\{01,23,49,57,68\},&
Q_1&=\{01,27,34,58,69\},\nonumber\\
Q_2&=\{04,12,38,57,69\},&
Q_3&=\{04,16,23,58,79\},\label{eq:p-matchings}\\
Q_4&=\{05,12,34,68,79\},&
Q_5&=\{05,16,27,38,49\}.\nonumber
\end{align}
Indeed, branch on the matching edge incident with vertex $0$.  Each of
$01,04,05$ leaves precisely the two completions displayed above.  Direct
comparison gives
\begin{equation}\label{eq:intersection}
|Q_j\cap Q_k|=1\quad(j\neq k),
\qquad |Q_j|=5.
\end{equation}

The table also shows that every edge of $P$ belongs to exactly two of the six
matchings.  For $j=0,\ldots,5$, define the lift $\widetilde Q_j$ by using the
image in $H$ of every edge of $Q_j$ and, at each replacement triangle, the
internal edge joining the two ports not incident with the selected external
edge.  Since each Petersen edge occurs in two of the $Q_j$, and each internal
triangle edge is selected by the two matchings using the opposite external
port, the six perfect matchings
\[
\widetilde Q_0,\ldots,\widetilde Q_5
\]
double-cover every edge of $H$.

\begin{lemma}[Unique matching cover]\label{lem:unique-cover}
After permuting the six indices,
\[
(N_0,\ldots,N_5)
=
(\widetilde Q_0,\ldots,\widetilde Q_5).
\]
\end{lemma}

\begin{proof}
Fix one replacement triangle.  A perfect matching uses at most one internal
edge there.  The three internal edges have total multiplicity six among the
six matchings $N_i$, because every edge belongs to exactly two of them.  Hence
each $N_i$ uses exactly one internal edge and therefore exactly one external
edge at this triangle.  This holds at all four triangles.  Contracting the
triangles maps every $N_i$ to a perfect matching of $P$, and the construction
above shows that its lift is unique.

Let $y_j$ be the multiplicity of $Q_j$ among the six contracted matchings.
Then $\sum_{j=0}^{5}y_j=6$.  For fixed $k$, sum the double-cover equations
over the five edges of $Q_k$.  By~\eqref{eq:intersection},
\begin{align*}
10
  &=\sum_{j=0}^{5}y_j|Q_k\cap Q_j|\\
  &=5y_k+\sum_{j\neq k}y_j\\
  &=5y_k+(6-y_k)\\
  &=6+4y_k.
\end{align*}
Thus $y_k=1$ for every $k=0,\ldots,5$.
\end{proof}

\begin{corollary}\label{cor:unique-6c4c}
Up to a permutation of its six members, $H$ has a unique unoriented
$6$-cycle $4$-cover, namely
\[
H-\widetilde Q_0,\ldots,H-\widetilde Q_5.
\]
\end{corollary}

\begin{proof}
Existence follows because the six lifted matchings double-cover every edge,
so their complements cover every edge four times.  Uniqueness follows from
Lemmas~\ref{lem:saturation} and~\ref{lem:unique-cover}.
\end{proof}

\section{The directional obstruction}

Use the vertex order
\begin{equation}\label{eq:vertex-order}
\begin{split}
p_{01},p_{04},p_{05},p_{10},p_{12},p_{16},u_2,u_3,
p_{40},p_{43},p_{49},{}\\
p_{50},p_{57},p_{58},u_6,u_7,u_8,u_9
\end{split}
\end{equation}
and direct every edge from the earlier endpoint to the later endpoint as the
reference direction.

For each $j$, the complement $H-\widetilde Q_j$ is the disjoint union of an
$11$-circuit $D_j$ and a $7$-circuit $D'_j$.  The component $D_j$ containing
$p_{01}$ has the following cyclic order:
\begin{align*}
D_0={}&(p_{01},p_{04},p_{40},p_{49},p_{43},u_3,u_8,
          p_{58},p_{57},p_{50},p_{05}),\\
D_1={}&(p_{01},p_{04},p_{40},p_{43},p_{49},u_9,u_7,
          p_{57},p_{58},p_{50},p_{05}),\\
D_2={}&(p_{01},p_{04},p_{05},p_{50},p_{57},p_{58},u_8,
          u_6,p_{16},p_{12},p_{10}),\\
D_3={}&(p_{01},p_{04},p_{05},p_{50},p_{58},p_{57},u_7,
          u_2,p_{12},p_{16},p_{10}),\\
D_4={}&(p_{01},p_{05},p_{04},p_{40},p_{43},p_{49},u_9,
          u_6,p_{16},p_{12},p_{10}),\\
D_5={}&(p_{01},p_{05},p_{04},p_{40},p_{49},p_{43},u_3,
          u_2,p_{12},p_{16},p_{10}).
\end{align*}
The six circuits $D'_j$ are listed in Appendix~\ref{app:components}.

Consider the five reference-directed edges
\begin{equation}\label{eq:selected-edges}
e_1=p_{01}p_{04},\quad
e_2=p_{05}p_{50},\quad
e_3=p_{12}p_{16},\quad
e_4=p_{43}p_{49},\quad
e_5=p_{57}p_{58}.
\end{equation}
A direct check shows that, for every $r\in\{1,\ldots,5\}$ and
$j\in\{0,\ldots,5\}$, either $e_r\in\widetilde Q_j$ or
$e_r\in E(D_j)$.  In particular, none of the five selected edges belongs to
$D'_j$.  Consequently, the orientation variables of the six $7$-circuits do
not occur in the five balance equations below.

Let $z_j\in\{-1,1\}$ record whether the orientation on $D_j$ agrees with the
displayed cyclic order.  Reading the directions of the five selected edges in
the six displayed circuits, equation~\eqref{eq:balance} gives
\begin{equation}\label{eq:matrix}
\begin{pmatrix}
 1& 1& 1& 1& 0& 0\\
-1&-1& 1& 1& 0& 0\\
 0& 0&-1& 1&-1& 1\\
-1& 1& 0& 0& 1&-1\\
-1& 1& 1&-1& 0& 0
\end{pmatrix}
\begin{pmatrix}
z_0\\
z_1\\
z_2\\
z_3\\
z_4\\
z_5
\end{pmatrix}
=0.
\end{equation}

\begin{lemma}[Five-row obstruction]\label{lem:obstruction}
The system~\eqref{eq:matrix} has no solution in $\{-1,1\}^6$.
\end{lemma}

\begin{proof}
Let $r_1,\ldots,r_5$ be the rows of the matrix in~\eqref{eq:matrix}.  Then
\[
r_1-r_2-r_3-r_4-r_5=(4,0,0,0,0,0).
\]
Left-multiplying~\eqref{eq:matrix} by $(1,-1,-1,-1,-1)$ therefore gives
$4z_0=0$, contradicting $z_0\in\{-1,1\}$.
\end{proof}

\begin{proof}[Proof of Theorem~\ref{thm:main}]
By Proposition~\ref{prop:eligible}, $H$ is a finite simple cubic
$3$-edge-connected graph.  If it admitted an oriented $6$-cycle $4$-cover,
Lemma~\ref{lem:saturation} would turn the six complementary edge sets into a
perfect-matching double cover.  By Lemma~\ref{lem:unique-cover}, after
permuting the indices these matchings would be
$\widetilde Q_0,\ldots,\widetilde Q_5$.  The orientations of the corresponding
$2$-factors would then produce signs satisfying~\eqref{eq:matrix}, contrary
to Lemma~\ref{lem:obstruction}.
\end{proof}

\section{Independent computational verification}

The proof above is self-contained and does not use computation as a premise.
For reproducibility, two separately written standard-library Python
implementations verify the finite data.  The first constructs $H$ from the
Petersen triangle expansion and derives the orientation equations from its
circuit decomposition.  The second starts from a separately numbered edge
list and independently enumerates the same structures.  The source code,
exact input files, expected output, and reproduction instructions are
packaged in the fixed release
\texttt{o6c4c-counterexample-v1.0.0.zip}, whose SHA-256 digest is
\begin{center}
\small\texttt{2243f4e96bc58f28c386a5b57f8af5a63ad43ebbf9efa323627f15be757b9d71}.
\end{center}
The fixed release is publicly available at
\url{https://github.com/Miracleofwinds/o6c4c-counterexample/releases/tag/v1.0.0}.

The verification includes:
\begin{itemize}[leftmargin=*,itemsep=0.2em]
\item connectivity after all $1+27+\binom{27}{2}=379$ deletions of at most two
      edges;
\item the six Petersen perfect matchings and all $18$ perfect matchings of $H$;
\item the unique double cover among $100{,}947$ multisets of six perfect
      matchings;
\item all twelve circuit components and all signs in~\eqref{eq:matrix}; and
\item infeasibility of all $2^{12}=4{,}096$ component-orientation assignments.
\end{itemize}

\section{Scope of the counterexample}

Because $H$ is cubic, every cycle in $H$ is a disjoint union of circuits.
Thus the proof uses only the componentwise cyclic orientation of ordinary
circuits; no convention concerning orientations at vertices of degree greater
than two is involved.

Disconnected cycles and repetitions among the six indexed members are
permitted throughout.  Consequently, the nonexistence result also applies to
any stronger formulation requiring the six underlying edge sets to be
distinct.

The graph $H$ does admit an unoriented $6$-cycle $4$-cover, uniquely up to
permutation, by Corollary~\ref{cor:unique-6c4c}.  The obstruction is therefore
genuinely orientational.  Finally, $H$ contains triangles and has cyclic
edge-connectivity $3$, so it lies outside the class of snarks of girth at
least $5$ and cyclic edge-connectivity at least $4$ tested computationally
in~\cite[Section~2.3]{Ulyanov}.

\appendix
\section{Exact numbered graph data}\label{app:data}

For direct reproduction, number the vertices in the
order~\eqref{eq:vertex-order}:
\[
\begin{array}{rcl@{\qquad}rcl@{\qquad}rcl}
0&:&p_{01} & 1&:&p_{04} & 2&:&p_{05}\\
3&:&p_{10} & 4&:&p_{12} & 5&:&p_{16}\\
6&:&u_2     & 7&:&u_3     & 8&:&p_{40}\\
9&:&p_{43} &10&:&p_{49}  &11&:&p_{50}\\
12&:&p_{57}&13&:&p_{58}  &14&:&u_6\\
15&:&u_7   &16&:&u_8     &17&:&u_9
\end{array}
\]
The $27$ edges, each listed in its reference direction, are
\begin{align*}
&(0,3),(1,8),(2,11),(4,6),(5,14),(6,7),(6,15),(7,9),(7,16),\\
&(10,17),(12,15),(13,16),(14,16),(14,17),(15,17),\\
&(0,1),(0,2),(1,2),(3,4),(3,5),(4,5),\\
&(8,9),(8,10),(9,10),(11,12),(11,13),(12,13).
\end{align*}
The first fifteen edges correspond to the Petersen edges in
\eqref{eq:petersen-edges}; the remaining twelve are the edges of the four
replacement triangles.

\section{Complete circuit decompositions}\label{app:components}

For completeness, the second component of $H-\widetilde Q_j$ is the following
$7$-circuit:
\begin{align*}
D'_0={}&(p_{10},p_{12},u_2,u_7,u_9,u_6,p_{16}),\\
D'_1={}&(p_{10},p_{12},u_2,u_3,u_8,u_6,p_{16}),\\
D'_2={}&(u_2,u_7,u_9,p_{49},p_{40},p_{43},u_3),\\
D'_3={}&(u_3,p_{43},p_{40},p_{49},u_9,u_6,u_8),\\
D'_4={}&(u_2,u_3,u_8,p_{58},p_{50},p_{57},u_7),\\
D'_5={}&(p_{50},p_{58},u_8,u_6,u_9,u_7,p_{57}).
\end{align*}
Together with the six $D_j$ in Section~5, these are all twelve components of
the six complementary $2$-factors.

\section{A minimal manual audit of the matrix}

To check~\eqref{eq:matrix} without a program, use the five edges in
\eqref{eq:selected-edges}.  For a fixed row and column, enter $0$ if the edge
belongs to $\widetilde Q_j$, $+1$ if the cyclic sequence for $D_j$ traverses
it in the reference direction, and $-1$ otherwise.  This gives the displayed
matrix entry by entry.  Its row combination is
\[
(1,-1,-1,-1,-1)
\begin{pmatrix}
 1& 1& 1& 1& 0& 0\\
-1&-1& 1& 1& 0& 0\\
 0& 0&-1& 1&-1& 1\\
-1& 1& 0& 0& 1&-1\\
-1& 1& 1&-1& 0& 0
\end{pmatrix}
=(4,0,0,0,0,0),
\]
which is the complete final obstruction.

\section*{Acknowledgements}

The counterexample was discovered with the assistance of Semidirect, an
automated mathematical research system.  All mathematical claims in this
paper were subsequently reduced to the explicit finite proof above and
checked by separately written verification implementations.

\begin{thebibliography}{9}
\bibitem{Fulkerson}
D.~R.~Fulkerson,
Blocking and anti-blocking pairs of polyhedra,
\emph{Math. Programming} \textbf{1} (1971), 168--194.

\bibitem{Ulyanov}
N.~Ulyanov,
Graph Puzzles I.1: Oriented Berge--Fulkerson Conjecture,
arXiv:2501.05348v3 [math.CO], 2026.

\bibitem{Zhang}
C.-Q.~Zhang,
\emph{Circuit Double Cover of Graphs},
London Mathematical Society Lecture Note Series 399,
Cambridge University Press, 2012.
\end{thebibliography}

\end{document}
